\chapter{Derivations for the Mental-Rotation Examples}
\label{annex:p19-mental-rotation}

\begin{annexitem}{annex:p19.rotation.fit.identification}{thm:p19-rotation-fit-and-identification}{Affine arithmetic and a limit of process identification}
\begin{proof}
Write \(x_i=\delta_i-\bar\delta\), \(y_i=t_i-\bar t\).
Using \(\sum x_i=\sum y_i=0\), expansion gives
\[
 \sum_i(t_i-a-b\delta_i)^2
 =\sum_i(y_i-bx_i)^2+
                  m(\bar t-a-b\bar\delta)^2.
\]
For fixed \(b\), the second term is uniquely minimized by
\(a=\bar t-b\bar\delta\). The first is
\(\sum y_i^2-2b\sum x_iy_i+b^2S_{\delta\delta}\).
Completing the square with \(S_{\delta\delta}>0\) proves
the unique fitted coefficients.
For the fictional table, centered angles are
\(-90,-45,0,45,90\) and centered times
\(-1428,-708,22,682,1432\).
Their square and cross-product sums are 20250 and 319950.
Division yields \(15.8\), and \(2628-15.8(90)=1206\).
Substitution gives the fitted values and residuals.
The squared residuals sum to \(36+9+484+841+100=1470\).
A rate \(v\) degrees/second contributes \(1000\delta/v\)
milliseconds, proving the reciprocal under its stated
process-assignment premise.

For nonidentification fix a finite grid \(\delta=j\Delta\),
\(a,b,\Delta>0\), and a common noise variable \(E\) supported
above \(-a\). One model has a continuously evolving operation
lasting \(b\delta\); a second has exactly \(j\) discrete updates,
each lasting \(b\Delta\). Add \(a+E\) in each model.
Their conditional time laws are identical, those of
\(a+bj\Delta+E\), despite different stipulated internal processes.
This proves nonidentification from those observations alone.
It makes no assertion about predictions for the remaining
behavioral or eye-movement observations.
\end{proof}
\end{annexitem}

\begin{annexitem}{annex:p19.feature.swap.benchmark}{thm:p19-feature-swap-benchmark}{A stipulated detection model}
\begin{proof}
Among \(\binom N2\) equally likely swaps, exactly
\(\binom{N-k}{2}\) avoid every remembered part.
Subtracting that fraction gives \(d_k\).
Total probability gives \(h=d_k+(1-d_k)g\), and the unchanged
trial's false-alarm probability is \(f=g\).
If \(g<1\), rearranging \(h=f+(1-f)d_k\) gives the correction.
If \(g=1\), every response is different and the division is
undefined; it supplies no detection estimate.
For \(N=4\), the undetected-swap counts are \(6,3,1,0,0\);
division by six gives the table.
For \(N=2\), the single possible swap involves a remembered
part both at \(k=1\) and at \(k=2\), proving the ceiling ambiguity.
Imperfect memory, dependence, biased guesses or a different
response rule would alter the model. The proof does not
assume these ideal premises true of the study participants.
\end{proof}
\end{annexitem}
